% BANKS-SOURCE: pdf=79 print=69 kind=section id=6.4
\subsection{The magnetic moment of a weakly coupled charged particle}

% BANKS-SOURCE: pdf=79 print=69 kind=prose
One of the most interesting calculations we can do with the external field method is
of the correction to the magnetic moment of a particle whose strongest interaction is

% BANKS-SOURCE: pdf=80 print=70 kind=figure id=6.4
\begin{figure}[htbp]
  \centering
  % Native TikZ reconstruction of Figure 6.4.
\begin{tikzpicture}[x=1cm,y=1cm,>=Stealth,baseline=(current bounding box.center)]
  \tikzset{BanksArrowLine/.style={line width=.65pt,postaction={decorate},
      decoration={markings,mark=at position .995 with {\arrow{Stealth[length=6pt]}}}},
    BanksPhoton/.style={line width=.65pt,decorate,
      decoration={snake,amplitude=4.5pt,segment length=13pt}}}
  \coordinate (v)  at (0,0);            % external-photon vertex
  \coordinate (au) at (-1.49,0.91);     % upper internal-photon attachment
  \coordinate (al) at (-1.44,-0.88);    % lower internal-photon attachment
  \coordinate (pu) at (-2.87,1.75);     % outgoing p+q
  \coordinate (pl) at (-2.90,-1.79);    % incoming p
  % outgoing fermion leg (arrowhead sits at the photon attachment; see E-033 —
  % the source draws this arrow pointing into the vertex, breaking the
  % continuity of the fermion line, and the adopted form reverses it)
  \draw[BanksArrowLine] (au) -- (pu);
  \draw[line width=.65pt] (au) -- (v);
  % incoming fermion leg
  \draw[BanksArrowLine] (pl) -- (al);
  \draw[line width=.65pt] (al) -- (v);
  % internal photon
  \draw[BanksPhoton] (au) -- (al);
  % external photon
  \draw[BanksPhoton] (v) -- (2.2,0);
  \fill (v) circle (2.4pt);
  \node[left] at (pu) {$p+q$};
  \node[left] at (pl) {$p$};
  \node at (-0.84,1.07) {$k+q$};
  \node at (-0.90,-0.78) {$k$};
  \node at (-2.32,0.05) {$p-k$};
  \node[right] at (2.28,0) {$q$};
\end{tikzpicture}

  \caption{One-loop vertex correction.}
  \label{fig:6.4}
\end{figure}

% BANKS-SOURCE: pdf=80 print=70 kind=prose
electromagnetic. The classic result of Dirac is that the gyromagnetic ratio of a point
particle of spin $\frac12$ is 2. What we will show is that purely electromagnetic interactions
correct that result. In fact, higher-order calculations and modern experimental tech-
niques have combined to make this calculation/measurement for the electron and the
muon one of the most precise agreements we have between theoretical physics and the
real world. The current state of the art is sensitive to the effects of virtual strongly
interacting particles and of the weak interactions. It can even be used to put inter-
esting bounds on physics at quite high energy scales. The electron and muon dipole
moments are known to 13 and 10 significant digits experimentally. For the electron
there is complete agreement between theory and experiment at that level of accuracy.
Recent (November 2006) results indicate a possibly significant deviation for the muon,
which might be an indication of physics beyond the standard model. We will calculate
only the leading correction to the anomalous magnetic moment. A general reference
for precision QED results is~\cite{banks-ref-34}.

We will employ the external field methods we have just learned. In the $D=4$
specialization used for the magnetic moment, take a field of the form
$F_{ij}=\epsilon_{ijk}B^k$, where $B^k$ is the constant value of the magnetic field,
in the rest frame of the particle. The electric field is assumed zero in this frame. We
will do the computation to leading order in the corrections to the result for the Dirac
equation in an external field. The reader should first do Problem 6.5, which shows that
the prediction of the Dirac theory is a gyromagnetic ratio of 2.

To first order in the external field, the invariant amplitude is given by
\begin{equation*}
  (-\ii e_0)\int \dd^Dq\,A^\mu(q)\bar u(p+q)\Gamma_\mu(p,q)u(p).
\end{equation*}
We have suppressed the spin indices on the incoming and outgoing spinors. We will
evaluate the 1PI vertex $\Gamma_\mu$ to one-loop order (Figure~\ref{fig:6.4}), taking into account that
to this order the bare mass parameter $m_0$ is equal to the physical mass $m$. Thus, the
diagram is given by

% BANKS-SOURCE: pdf=80 print=70 kind=display id=6.4a
\begin{equation*}
  \begin{aligned}
  \bar u\Gamma_\mu u
  ={}&\frac{(-\ii)(-\ii e_0)^2}{(2\pi)^D}\bar u(p+q)\int \dd^Dk\,
  \frac{\eta_{\lambda\kappa}+(k-p)_\lambda(k-p)_\kappa/\mu^2}
  {(k-p)^2+\mu^2-\ii\epsilon}\\[-2pt]
  &\quad\times
  \frac{\gamma^\lambda[(\slashed k+\slashed q+m)\gamma_\mu(\slashed k+m)]\gamma^\kappa}
  {[(k+q)^2+m^2-\ii\epsilon][k^2+m^2-\ii\epsilon]}u(p).
  \end{aligned}
\end{equation*}

% BANKS-SOURCE: pdf=81 print=71 kind=prose
Let us first deal with the term which apparently diverges as $\mu\to0$. Using the
identities $\bar u(p+q)(\slashed p-\slashed k)=\bar u(p+q)(m-\slashed k-\slashed q)$ and
$(\slashed p-\slashed k)u(p)=(m-\slashed k)u(p)$,
we see that the denominators in the two fermion propagators are canceled out by the
numerators, and we are left with a term proportional to
\begin{equation*}
  \frac{e_0^2}{\mu^2}\bar u\gamma^\mu u\int\frac{\dd^Dk}{k^2+\mu^2-\ii\epsilon}.
\end{equation*}
This term can be absorbed into a rescaling of the fermion fields because it has the
same form as the tree-level vertex. The precise value of this renormalization depends
on the method we use for cutting off ultraviolet divergences in the theory, a subject to
which we will return in Chapter 9. There we will distinguish between cut-off procedures
that preserve gauge invariance and those which don’t. With a gauge-invariant regula-
tor, like the analytic continuation of Feynman diagrams in the space-time dimension\footnote[5]{This is called \emph{dimensional regularization} and abbreviated DR.}
the apparently divergent integral in the previous expression is actually finite and pro-
portional to $\mu^2$. We have to be careful about the regulation procedure because, as we
have noted, massive electrodynamics is really a gauge theory in its Higgs phase. The
upshot of this discussion is that we can ignore the longitudinal term in the gauge-boson
propagator, and replace it by
\begin{equation*}
  -\ii\frac{\eta_{\mu\nu}}{p^2+\mu^2-\ii\epsilon}
\end{equation*}
in this diagram.  Its massless Feynman-gauge limit is
\(-\ii\eta_{\mu\nu}/(p^2-\ii\epsilon)\).

Before continuing the computation let us stop to ask what we expect to get. The
current matrix element is $\bar u(p+q)\Gamma_\mu u(p)$, and the vertex function $\Gamma_\mu$ can be expanded
in terms of the linearly independent Dirac matrices $1$, $\gamma_5$, $\gamma_\mu$, $\gamma_\mu\gamma_5$, $\gamma_{\mu\nu}$. The two expres-
sions with $\gamma_5$, as well as anything else involving the Levi-Civita symbol, cannot appear,
because massive electrodynamics conserves parity. We must combine the matrices with
the 4-vectors $p,q$ to make a 4-vector, and, when we sandwich $\Gamma_\mu$ between on-shell
spinors, we can use
\begin{equation*}
  (\slashed p-m)u(p)=(\slashed p+\slashed q-m)u(p+q)=0
\end{equation*}
and the conjugates of these equations to prove that
\begin{equation*}
  \bar u(p+q)\slashed q\,u(p)=0
\end{equation*}
and
\begin{equation*}
 \bar u(p+q)(2p^\mu+q^\mu)u(p)
  =\bar u(p+q)[\gamma^{\mu\nu}q_\nu-2m\gamma^\mu]u(p).
\end{equation*}
The last identity, named after Gordon, was proved in the exercises to Chapter 5. Using
these identities, it is easy to see that anything in $\Gamma_\mu$ can be reduced to
\begin{equation*}
  F_1(q^2,m^2)\gamma_\mu+F_2(q^2,m^2)\frac{\gamma_{\mu\nu}q^\nu}{2m}
  +F_3(q^2,m^2)q_\mu,
\end{equation*}

% BANKS-SOURCE: pdf=82 print=72 kind=prose
when evaluated between on-shell spinors. The form factors $F_i$ depend on the Lorentz
invariants $q^2$, $p^2=-m^2$ and $2p\!\cdot q=-q^2$.

\BanksImplicitHook{I-CH06-005}

The current matrix element is conserved, which means $q^\mu\bar u\Gamma_\mu u=0$. This is satisfied
for arbitrary $F_{1,2}$ and implies that $F_3=0$. In our calculation, we will seek to reduce
all expressions in $\Gamma_\mu$ (after loop integration) to linear combinations of $(2p+q)^\mu$, $\gamma^\mu$,
and $q^\mu$. The diligent reader will prove that, once this is done, the coefficient of $q^\mu$
vanishes. The Gordon identity will allow us to identify the individual $F_{1,2}$. An extremely
important point of the analysis is that, for $q^2=0$, which is the only thing our constant
external field probes, $F_1$ serves merely to \emph{renormalize} the value of the electric coupling
$e_0$. Thus, it cannot change the gyromagnetic ratio of 2 implied by the Dirac equation
for any coupling. The change in the gyromagnetic ratio is due entirely to the second
form factor. $F_1$ is called the electric and $F_2$ the magnetic form factor.

\BanksImplicitHook{I-CH06-006}

At this point in our studies, this is fortunate. A direct calculation of $F_1(0)$ (though
not its derivatives w.r.t. $q^2$) would lead to infinite integrals, whose meaning I am not
yet prepared to explain. We will find, however, that $F_2$ is completely finite and
unambiguous. This can be seen by dimensional analysis. The integral defining $\Gamma_\mu$ contains
four powers of loop momentum in the denominator. It can, at most, give rise to a
logarithmic divergence. If we differentiate it w.r.t. $q$, we get a finite integral. Only the
first term in the Taylor expansion of $\Gamma_\mu$ diverges. The term containing $F_2$ vanishes at
$q=0$ and will therefore be finite. This sort of dimensional analysis will be the key to
the understanding of renormalization, which we will achieve in Chapter 9.

For those readers who wish to practice their skills by calculating the entire vertex
function, I will introduce a method to render all integrals finite. It consists in changing
the denominator in the photon propagator to
\begin{equation*}
  (p^2+\mu^2-\ii\epsilon)K(p^2/\Lambda^2),
\end{equation*}
where $K$ is a smooth function, which is 1 for $|p^2|<\Lambda^2$ and blows up at least exponentially rapidly
above that. We will argue in Chapter 9 that, as long as $\Lambda$ is much larger than the electron
and photon masses and the external momenta, the answers for physical amplitudes will be
independent of the detailed shape of $K$.

I will also employ an identity, invented by Feynman (its derivation follows in the
problems), which states that
\begin{equation*}
  \frac{1}{A_1A_2A_3}=\int_0^1 dx\,dy\,dz\,\delta(1-x-y-z)\frac{2}{\Delta^3},
\end{equation*}
where
\begin{equation*}
  \Delta=xA_1+yA_2+zA_3.
\end{equation*}
For our purposes, the three $A_i$ are the Feynman denominators, and
\begin{equation*}
  \Delta=k^2+2k\!\cdot(yq-zp)+yq^2+zp^2+(x+y)m^2+z\mu^2-\ii\epsilon.
\end{equation*}
If we introduce $r=k+yq-zp$, then
\begin{equation*}
  \Delta=r^2+U-\ii\epsilon,
\end{equation*}
with
\begin{equation*}
  U=xyq^2+(1-z)^2m^2+z\mu^2.
\end{equation*}

% BANKS-SOURCE: pdf=83 print=73 kind=prose
This is extremely useful, because the numerator contains terms of zeroth, first, and
second order in $r_\mu$. The linear terms vanish upon integration, by virtue of Lorentz
invariance, as long as we are careful to regulate the divergences with a Lorentz-invariant
prescription like DR, or our modified photon propagator. The quadratic terms are also
simplified by Lorentz invariance:
\begin{equation*}
  \int \dd^D r\,\frac{r_\mu r_\nu}{\Delta^3}
  =\frac{\eta_{\mu\nu}}{D}\int \dd^D r
    \left(\frac1{\Delta^2}-\frac{U}{\Delta^3}\right).
\end{equation*}

Using the anti-commutation relations, we can rewrite the numerator of the diagram
as
\begin{equation*}
  \mathrm{Num}=\gamma^\lambda[\slashed q\gamma_\mu(\slashed k+m)+(k^2+m^2)\gamma_\mu
  -2k_\mu(\slashed k+m)]\gamma_\lambda.
\end{equation*}
We use the contraction identities from Appendix C,
\begin{equation*}
  \gamma_\lambda\gamma_\mu\gamma^\lambda=2\gamma_\mu,
  \qquad
  \gamma^\lambda\gamma_\mu\gamma_\nu\gamma_\lambda=4\eta_{\mu\nu},
  \qquad
  \gamma^\lambda\gamma_\mu\gamma_\nu\gamma_\alpha\gamma_\lambda
  =-4\eta_{\nu\alpha}\gamma_\mu-2\gamma_\nu\gamma_\alpha\gamma_\mu,
\end{equation*}
and the fact that $\bar u(p+q)\slashed q\,u(p)=0$, to rewrite this as
\begin{equation*}
  \mathrm{Num}=4mq_\mu+4k_\mu(2m-\slashed k)
  +2\gamma_\mu[k^2+m^2-\slashed k\slashed q].
\end{equation*}

We substitute $k=r+zp-yq$ and use the relations for integrals over $r$ to express
this, between the on-shell spinors, as
\begin{equation*}
  \begin{aligned}
  \mathrm{Num}={}&\gamma_\mu\left[\Delta-U+2m^2(1-2z-z^2)
       -2(1-y)(y+z)q^2\right]\\
  &+(2p+q)_\mu\,2mz(1-z)
    +q_\mu\,2m(2-3z+z^2-4y+2yz).
  \end{aligned}
\end{equation*}
We have again used the fact that the numerator is sandwiched between on-shell spinors.
Note in particular that
\begin{equation*}
\bar u(p+q)\gamma_\mu\slashed q\,u(p)
  =-2\bar u(p+q)(p_\mu+q_\mu+\gamma_\mu m)u(p).
\end{equation*}

We have also reintroduced $x=1-y-z$, in order to make it clear that the term
proportional to $q_\mu$ is odd under interchange of $x$ and $y$. It therefore vanishes when
integrated. We now use the Gordon identity to find that the magnetic form factor is
given by

% BANKS-SOURCE: pdf=83 print=73 kind=equation id=6.17
\begin{equation}
  -\frac{\ii e^2}{(2\pi)^D}\bar u(p+q)\frac{\gamma^{\mu\nu}q_\nu}{2m}u(p)
  \int \dd^D r\int_0^1dx\,dy\,dz\,\delta(1-x-y-z)
  \frac{m^2z(1-z)}{\Delta^3}.
  \label{eq:6.17}
\end{equation}
Note that we have left off the effect of the regulator for the photon propagator, as
well as the photon mass. In fact, we will find that the answer for the magnetic form
factor is completely finite. In Chapter 9 we will learn the reason for the UV finiteness:
expressed as corrections to the quantum Lagrangian for the Dirac field, the terms in the
Taylor expansion of the magnetic form factor around $q=0$ are all operators of dimen-
sion higher than four. In a theory whose couplings have non-negative mass dimension

% BANKS-SOURCE: pdf=84 print=74 kind=prose
(which we will learn to call a renormalizable theory at the Gaussian fixed point), such
terms are independent of the cut-off.

The $r$ integral is done by analytic continuation to Euclidean space, following the $\ii\epsilon$
prescription. In general $D$ dimensions,
\begin{equation*}
  \int\frac{\dd^D r_E}{(r_E^2+U)^3}
  =\frac{\pi^{D/2}}{2}\,\Gamma\!\left(3-\frac D2\right)U^{D/2-3}.
\end{equation*}
For the $D=4$ specialization used below this is
\begin{equation*}
  \int\frac{\dd^4r_E}{(r_E^2+U)^3}
  =\frac1U(2\pi^2)\int_0^\infty\frac{t^3\,\dd t}{(t^2+1)^3}.
\end{equation*}
The factor $2\pi^2$ is the 3-volume of the unit sphere embedded in four dimensions. The
Lorentzian signature integral is $+\ii$ times the Euclidean one. The factor of $\ii$ comes from
$r^0=\ii r_E^4$, with $r^2+U=r_E^2+U$. The result is

% BANKS-SOURCE: pdf=84 print=74 kind=equation id=6.18
\begin{equation}
  \frac{e^2}{4\pi^2}\bar u(p+q)\frac{\gamma^{\mu\nu}q_\nu}{2m}u(p)
  \int_0^1dx\,dy\,dz\,\delta(1-x-y-z)
  \frac{m^2z(1-z)}{(1-z)^2m^2+xyq^2}.
  \label{eq:6.18}
\end{equation}
At this point, we have set $\mu^2=0$, since we will obtain a finite result in this limit. The
electric form factor has an infrared divergence when $\mu=0$, reflecting the zero prob-
ability for scattering without any emission of massless photons, but, to this order in
perturbation theory, the magnetic form factor is IR finite. For $q^2=0$ it is easy to do the
remaining integrals. They give $F_2(0)=\alpha/(2\pi)$, where $\alpha=e^2/(4\pi)$ is the fine-structure
constant. In the exercises, you will verify that $F_2(0)=(g-2)/2$, also called the anoma-
lous magnetic momentum of the electron ($g$ is the gyromagnetic ratio). This result was
first obtained by Schwinger in 1948. Since then, experimental and theoretical determi-
nations have competed with each other in precision, with no definite discrepancy yet
having been found for the electron magnetic moment. Possible deviations for the muon
moment, if they exist, are probably indications for physics beyond the standard model.
We will learn more about the standard model of particle physics in Chapter 8.

\BanksImplicitHook{I-CH06-007}
